% EGA V, first essay, §13.4. Recollated against actual French physical80–82
% and English physical52–53. Original order a,b,d,c,e is recorded below.
\egas{13.4}{注}
\source{法稿物理80页中段至82页末；英编稿物理52页末至53页末}

\noindent a) 我们已经指出，对13.2的成立来说，$k$为完美域的假设是必要的。相反，“$k$无限”这一假设很可能是多余的。我们没有尝试为$k$有限的情形给出专门的论证，只指出：在这种情形下，对$k$的代数闭包应用13.2，再用通常的论证，便可找到$k$的一个有限扩张$k'$，使得对于$X_{k'}$中位于$x$上方的点，都存在相对于$k'$的多重初等开邻域。\cmt{V21-REM001}{法稿80页明确写$X_{k'}$及复数“les points”，并以手插“pour”引出这些点；英文转写的概形字母和单数不能据此制造唯一点断言。旧注D17-M081所作$X_{k'}$订正，现可直接由法稿确认。前句“很可能”仍是作者的推测，不能改写为本稿已经证明了有限域上无需扩域的定理。有限扩张下降的范围在本版报告中另行说明。}

\noindent b) 如果在13.2中放弃$X$几何不可约这一假设，结论当然不能原样成立，因为一个多重初等代数概形是几何不可约的。原稿提出的较弱形式，是在定义13.1中删去“连通”一词。不过，仅作这一修改还不足以涵盖零维情形；若要得到包括零维在内的推广，还须对零维起始对象另作约定。\cmt{V21-REM002}{原文后两句断言：“不过，结论仍以一种较弱的形式成立；在定义13.1中删去‘连通’一词，就得到这种形式。我们已经给出的证明表明这一点。”此项无条件断言有零维反例：取非平凡有限可分扩张$L/k$，如$\mathbb C/\mathbb R$，则$\Spec L$在$k$上光滑；删去连通后，每个非空初等态射仍为相对维数1的光滑满射，所以长度$r$的塔维数为$r$，而零次复合只给$\Spec k$。故$\Spec L\to\Spec k$不是这种塔。正文明确所缺的零维补约，未宣称任一修改后的定义已使所有高维推广成立。原定义的几何不可约性及反例详见报告V21-REM-E02。}

\noindent c) （或可收入13.2的陈述。）沿用13.2的记号，设$\Phi$是$X$的一个非空有限子集，由$X$的光滑点组成，并假设$\Phi$包含在$X$的某个仿射开集中。则$\Phi$具有一个由多重初等开邻域组成的基本邻域系。显然，可以假设$\Phi$由闭点组成。\cmt{V21-REM003}{原稿没有明说“非空”；若允许$\Phi=\varnothing$并要求在每个开邻域中取多重初等开集，取空开邻域就不能满足结论。正文明示实际应用所需的非空条件。法稿纸面顺序为a)、b)、d)、c)、e)，并带移排箭头；本版沿用英编稿及既有合订稿按编号a)--e)排列，所有段落均保留。换成闭点集须在给定的光滑仿射有限型邻域中选取闭特化；含这些特化的开集也含原点。有限型条件由光滑性取得，并非任意空间上的无条件结论。}

证明本质上相同，只须将13.3略作改写如下：存在$\mathbf P^r$中一个余维为$n$的线性子簇$C$，它不与$\Phi\cup Y$相交，并且对于每个$x_i\in\Phi(\bar k)$，由$C$与$x_i$张成的余维为$n-1$的线性子簇$L_i$具有13.3的性质a)、b)、c)。\cmt{V21-REM004}{法稿81页明确保留了$C$的余维$n$；$x_i$后的手改是$\Phi(\bar k)$，而非只取$k$-有理点。这里把闭点集赋予约化结构，$\Phi(\bar k)$表示它的几何点集。英文漏掉余维数及上横线，旧注D17-M083、084分别据上下文补数并保留闭点/有理点疑问；前一数值和后一几何点读法现在都有原稿依据。对非有理点取张成空间时在$\bar k$上进行，不能假称每个$L_i$本来都在$k$上定义。}

为验证这一点，注意只须验证一般的$C$具有上述性质。而对这样的$C$，每个$L_i$在经过$x_i$、余维为$n-1$的线性子簇$L$之中都是一般的；因此可以应用13.3的原有形式，或者至少应用我们已经证明的形式，即：每个经过$x$且充分一般的$L_0$都具有性质a)、b)、c)。\cmt{V21-REM005}{法稿手插“pour une telle”连接两句，末段“qui était toute”也可辨；英文中的连续问号和法语占位不再作为法稿缺文保留。法稿及英文末次比较$L$的余维仍写$n$，但同段$L_i$的余维是$n-1$，且$L_i=\langle C,x_i\rangle$；这里将该处订为$n-1$，原错记录于此。对有限个几何点同时选取$C$、再下降到$k$的步骤，在报告中作为译者补证展开，不能只检验$\Phi(k)$。}

\noindent d) 按照7.1所说明的方法，当把基域$k$换成一般的基预概形时，可以给出13.2的若干变式。我们指出下面一条（不作证明）：设$f\colon X\to S$是平坦射影态射，其几何纤维不可约并满足$(R_2)$；设$S'$是$X$的一个在$S$上有限的子预概形，$s\in S$，并假设对每个位于$s$上方的$x\in S'$，$X_s$在$x$处在$\kappa(s)$上光滑。则存在$s$的一个开邻域$U$，以及$X|_U$中$S'|_U$的一个开邻域$V$，使得$V\to U$是多重初等态射。\cmt{V21-REM006}{本段保留作者明说“不作证明”的草案身份。法稿81页写“几何纤维不可约且$(R_2)$”，所以$(R_2)$也施于几何纤维；旧注D17-M088关于是否为几何条件的疑问现予撤销。法稿的$s\in S$也清楚，英稿误作$x\in S$。态射靶字母在法稿右边缘不完整，本版据随后始终使用的$S$统一记基底，不能声称裁切字形已经读出。一般基底下，仅取Zariski开邻域是否足够、有限剩余域及有限表示条件如何处理，原段没有证明；本版仍将这些列为未决，未用7.1的引用冒充解决。}

若$Y$是$X$的一个不与$S'$相交的闭子预概形，且$Y$在$S$上不光滑的点所成的集合$Z$满足
\[
\dim Z_s\leq\dim X_s-2,
\]
则在上述结论中，还可以取$V\subset X-Y$。\cmt{V21-REM007}{法稿明确把非光滑点集记作$Z$，并以手插“satisfasse”连接维数不等式；“不与$S'$相交”也是清楚的手插条件。英文的点集字母$X$及残缺连接语不代表原法稿在此缺文。旧注D17-M090据英文留下的文字疑问现撤销；数学上对一般基底的未证范围不因此消失。}

\noindent e) 13.2值得注意的原因之一，是多重初等代数概形$U$具有特别简单的“拓扑”结构。\cmt{V21-REM008}{法稿81页把这里写成“初等代数概形”，随后却讨论任意维数$n$及自由群的逐次扩张；按13.2的对象及13.1的定义，正文明确为“多重初等”。原稿在该处有覆打及边缘标记，不把这个数学语义订正伪称为辨认出了完整的“poly”字样。}
例如，基域为复数域时，以$U^{\mathrm{an}}$表示与$U$关联的解析空间，则同伦群
\[
\pi_i(U^{\mathrm{an}})=0\qquad(i\geq2),
\]
而$\pi_1(U^{\mathrm{an}})$是自由群的逐次扩张。\cmt{V21-REM009}{法稿82页写的是$U^{\mathrm{an}}$，英文出现的$Y^{\mathrm{en}}$、$U^{\mathrm{en}}$均不保留为作者记号。法稿的指标是$i\ne1$；正文按正次数同伦群写成$i\geq2$，连通分支集合$\pi_0$另由连通性处理。原稿随后删去“又因$U$连通而$U^{\mathrm{an}}$连通”这一过渡，并改接“所以”；此为作者删改层，不把被删句补回现行正文。}

所以，$U^{\mathrm{an}}$是一个为$\pi_1$分类的“$K(\pi_1,1)$空间”。更确切地说，它的普遍覆盖空间同胚于$\mathbb C^n$，从而当然可缩；这个覆盖是一个结构群为$\pi_1$的“普遍主纤维丛”。\cmt{V21-REM010}{这里$n=\dim U$，并按13.1理解非空多重初等概形及非空有限平展边界。法稿$n$为上标，英文却排成$C_n$，正文据法稿恢复。原稿只说同胚，没有声称双全纯。本版报告使用保边界的Ehresmann平凡化、穿孔紧曲面的覆盖分类和对初等塔的归纳，分别核查高同伦消失、基本群扩张及较强的同胚断言；不能仅由普遍覆盖可缩便反推它同胚于$\mathbb C^n$。旧注D17-M093的“尚未核查”是历史状态。}
